\section{Comparison across shape-model versions}
\label{sec:crossversion}

Sections~\ref{sec:slopesens} and~\ref{sec:varrobust} varied the mesh while
holding the shape model fixed. This section varies the shape model itself. The
question is practical: most asteroids will never be visited, and their shapes
are known only from ground-based radar or lightcurve inversion. Which surface
gravity quantities, if any, survive that substitution?

\subsection{The controlled case: Itokawa}
\label{sec:itokawaxv}

Itokawa is the rare object for which two independent shape models of comparable
extent exist: the photogrammetric model built from Hayabusa imaging, with
$49\,152$ facets, and a radar-derived model with $3\,688$. Their bounding boxes
agree to $6.5\%$, $3.9\%$ and $13.3\%$ along the three principal axes, so the
two models describe a body of the same size and gross figure; what differs is
the detail they resolve. Both are evaluated at the same density and rotation
period, so the comparison isolates the effect of shape-model provenance.

Because Section~\ref{sec:slopesens} showed that resolution alone moves the
slope, the two models must be compared at matched resolution or the test is
meaningless. The radar model has $3\,688$ facets; the corresponding
spacecraft-model values of the bulk descriptors are obtained by linear
interpolation of the sequences of Tables~\ref{tab:slopeseries}
and~\ref{tab:varseries} between the $2\,183$- and $6\,404$-facet meshes. The
distribution statistics cannot be interpolated in the same way and are taken
from the nearest available direct runs, at $4\,857$ facets for the spacecraft
model against $3\,688$ for the radar model; the comparison of the tails is
therefore at approximately rather than exactly matched resolution, a caveat we
return to below. Table~\ref{tab:itokawaxv} gives the result, and
Figure~\ref{fig:crossversion}(a) shows the two slope distributions.

\begin{table}[t]
\centering
\small
\caption{Itokawa: spacecraft-derived and radar-derived shape models compared at
$\Nf \simeq 3\,688$, at $\dens = 1.95\ \mathrm{g\,cm^{-3}}$ and
$P = 12.132$~h. Spacecraft-model values for the bulk descriptors are
interpolated to the radar model's facet count; distribution statistics are from
direct runs at $4\,857$ and $3\,688$ facets respectively.}
\label{tab:itokawaxv}
\begin{tabular}{lrrr}
\toprule
Quantity & Spacecraft & Radar & Difference \\
\midrule
$\slopeavg$ (area-weighted)   & $14.87^{\circ}$ & $14.75^{\circ}$ & $-0.8\%$  \\
$\potvarn$                    & $0.08545$       & $0.07874$       & $-7.9\%$  \\
\midrule
Mean slope (unweighted)       & $15.91^{\circ}$ & $13.40^{\circ}$ & $-15.8\%$ \\
Maximum slope                 & $149.45^{\circ}$& $41.32^{\circ}$ & --- \\
Facets above $30^{\circ}$     & $4.8\%$         & $1.7\%$         & --- \\
Facets above $90^{\circ}$     & $2.29\%$        & $0.00\%$        & --- \\
\bottomrule
\end{tabular}
\end{table}

\begin{figure}[tp]
\centering
\includegraphics[width=\textwidth]{fig5.pdf}
\caption{Shape-model detail and the slope distribution. (a) Slope distribution
of Itokawa from the spacecraft-derived ($N_f = 4\,857$) and radar-derived
($N_f = 3\,688$) shape models. The bulk descriptors, compared at a common facet
count by interpolation, agree closely---their area-weighted mean slopes differ
by $0.8\%$---but the radar model contains no steep tail.
(b) Effect of Laplacian smoothing at fixed facet count. Ten iterations bring the
spacecraft model of Itokawa to $12.90^{\circ}$, within $3.7\%$ of the value
obtained independently from the radar model (dashed line), supporting the
interpretation that the difference between the two models is dominated by
facet-scale relief.}
\label{fig:crossversion}
\end{figure}

The two halves of that table point in opposite directions, and the contrast is
the result of this section.

The bulk descriptors, compared at the same facet count, agree well: the
area-weighted mean slope differs by less than one per cent and the normalized
geopotential dispersion by under eight per cent. A ground-based shape model, carrying an order of magnitude fewer facets
and no knowledge of the surface below the radar resolution, reproduces both
global summaries of the gravity field to within the accuracy with which the
assumed bulk density is usually known.

The distribution tails do not agree at all. The radar model contains no facet
steeper than $41^{\circ}$ and places $1.7\%$ of the surface above $30^{\circ}$;
the spacecraft model reaches $149^{\circ}$, places $4.8\%$ above $30^{\circ}$,
and assigns $2.29\%$ of its facets slopes exceeding $90^{\circ}$---overhanging
surfaces that the radar model does not represent at all. Nothing in the radar
model's global agreement gives any warning that its steep tail is absent.

The unweighted mean slope sits between the two behaviours, differing by
$15.8\%$ where the area-weighted mean differs by $0.8\%$. The reason is that
the spacecraft model's steep facets are also its smallest: they occupy a minor
fraction of the surface area but a substantial fraction of the facet count, so
they dominate an unweighted average and are correctly suppressed in an
area-weighted one. This is a further argument for the area weighting adopted in
Eq.~\eqref{eq:slopeavg}, and a caution that the two averages should not be used
interchangeably when models of different resolution are involved.

\subsection{Consistency with the smoothing experiment}
\label{sec:xvsmooth}

Section~\ref{sec:slopeorigin} attributed the resolution dependence of the slope
to relief that coarse meshes fail to resolve. If the difference between the two
Itokawa models is the same effect---the radar model simply never having
recorded that relief---then artificially removing the relief from the
spacecraft model should bring it toward the radar model's value.

It does. Ten Laplacian iterations applied to the spacecraft model
(Table~\ref{tab:smooth}) reduce its mean slope to $12.90^{\circ}$, within
$3.7\%$ of the $13.40^{\circ}$ obtained from the independent radar model, as
Figure~\ref{fig:crossversion}(b) shows. The
agreement is not exact and should not be over-read: smoothing and coarse
observation are different operations, and the residual difference is comparable
to what remains between the models generally. But that a purely internal
operation on one model reproduces, to within a few per cent, the value obtained
from an independently constructed model is a useful consistency check on the
interpretation.

\subsection{A cautionary counterexample: Kleopatra}
\label{sec:kleopatraxv}

Two shape models are also available for Kleopatra, and applying the same
comparison to them illustrates a prerequisite that is easy to overlook.

Compared at $\Nf \simeq 2\,292$, the two models disagree by $39.9\%$ in
area-weighted mean slope ($19.90^{\circ}$ against $11.96^{\circ}$) and by
$36.1\%$ in geopotential dispersion ($0.03683$ against $0.02354$). Taken at face
value this would overturn the Itokawa result.

It cannot be taken at face value. The two models do not describe the same body
at different levels of detail: their bounding boxes differ in scale. The
alternative model has a long axis of $276.3$~km against roughly $217$~km for
the radar model, while the intermediate and short axes agree to a few per cent,
so the two differ in aspect ratio ($3.6{:}1$ against $2.7{:}1$) as well as in
overall size. At fixed assumed density a longer body contains more mass, its
self-gravity is stronger relative to the fixed centrifugal term, and its slopes
are correspondingly lower---which is the direction of the observed
discrepancy. The $36$--$40\%$ differences therefore conflate scale, elongation
and surface detail, and cannot be attributed to shape-model provenance.

We report the comparison rather than discarding it, because the lesson is
useful. Before two shape models of one body are compared, their bounding boxes
should be checked against each other and against the published dimensions; this
is among the diagnostics listed in Section~\ref{sec:shapes} for exactly this
reason. Kleopatra fails that check and Itokawa passes it, which is why the
conclusions of this section rest on Itokawa alone.

\subsection{Implications for bodies known only from the ground}
\label{sec:xvimplications}

For the large majority of small bodies, a ground-based shape model is the only
one that will ever exist. The Itokawa comparison suggests a division of labour
for such objects.

In the one controlled case available to us, global descriptors of the surface
gravity field are largely preserved: the area-weighted mean slope is recovered
to within one per cent and the normalized geopotential dispersion to within
eight per cent, the latter comparable to the uncertainty contributed by the
assumed bulk density. Where
the goal is to characterise a population, to rank bodies by how far their
surfaces sit from an equipotential, or to provide first-order context for
mission planning, such models are adequate.

Anything that depends on the tail of the slope distribution is not recoverable.
Identifying the steepest terrain, estimating the fraction of surface above the
angle of repose, or locating overhangs requires relief that a ground-based
model does not contain, and the resulting figures will be systematically and
silently optimistic. A coarse model does not merely give a noisier answer to
these questions; it gives a confidently wrong one.

Demonstrated as it is for a single object, this comparison is not a general
calibration. It does, however, provide a first quantitative reference point for
the size of the discrepancy to be expected---close agreement in the bulk
descriptors, order-unity disagreement in the tails---and the same test can be
repeated wherever a second, independent shape model of a body becomes
available.
